\subsection*{習題 7.1：提示與參考解答}
\sourcepage{合成測試來源 S2，p.202}
\textbf{提示。} 其他條件不變時，先用式\eqref{ch7:rate}觀察厚度如何影響通量，再乘以面積；比值要與例題的 $200\,\mathrm W$ 比較。

\textbf{參考解答（S2 原題；以下為 AI 求解，來源未附答案）。}\par
已知 $L=0.040\,\mathrm m$、$k=0.80\,\mathrm{W\,m^{-1}\,K^{-1}}$、$A=0.50\,\mathrm{m^2}$，溫差仍為 $10\,\mathrm K$；沿用相同模型及由熱面指向冷面的正方向。由式\eqref{ch7:rate}，
\[
\text{(a)}\quad q_x=\frac{(0.80)(10)}{0.040}=200\,\mathrm{W\,m^{-2}}\quad(+x),
\]
\[
\text{(b)}\quad \dot Q=Aq_x=(0.50)(200)=100\,\mathrm W,\qquad
\text{(c)}\quad \frac{\dot Q_{\rm new}}{\dot Q_{\rm example}}=\frac{100}{200}=\frac12.
\]
厚度加倍而其他量不變，使梯度大小減半，故通量和總傳熱率都減半，與上述比值一致。通量是每單位面積的傳熱率，所以單位是 $\mathrm{W\,m^{-2}}$；總傳熱率乘入截面積後單位是 $\mathrm W$。
